% Generated draft: requires mathematical and source review.
\documentclass[11pt]{article}
\usepackage[margin=25mm]{geometry}
\usepackage{fontspec}
\setmainfont{Latin Modern Roman}
\usepackage{amsmath,amssymb,mathtools,unicode-math}
\setmathfont{Latin Modern Math}
\usepackage{enumitem,needspace,xurl,hyperref}
\setlength{\parindent}{0pt}
\setlength{\parskip}{5pt}
\setlength{\emergencystretch}{3em}
\newcommand{\N}{\mathbb{N}}
\newcommand{\Z}{\mathbb{Z}}
\newcommand{\Q}{\mathbb{Q}}
\newcommand{\R}{\mathbb{R}}
\newcommand{\C}{\mathbb{C}}
\newcommand{\F}{\mathbb{F}}
\newcommand{\A}{\mathbb{A}}
\newcommand{\PP}{\mathbb{P}}
\newcommand{\E}{\mathbb{E}}
\newcommand{\eps}{\varepsilon}
\newcommand{\dd}{\,\mathrm{d}}
\begin{document}
\section*{UnsolvedMath ID: 30004599}
\textbf{Source problem code:} catalogue entry
\subsection{Definitions and mathematical statement}
The principal question in the supplied cleaned record concerns further rank-three simplicial arrangements. The record does not specify the underlying field, the equivalence relation on arrangements, or the collection of already known examples. The following gives a conventional real, central interpretation, conditional on these choices being intended.

Let \(V\) be a finite-dimensional real vector space and let \(\mathcal A\) be a finite set of linear hyperplanes in \(V\). Set
\[
T=\bigcap_{H\in\mathcal A}H,
\qquad \operatorname{rank}(\mathcal A)=\dim(V/T).
\]
For rank three, write \(W=V/T\) and \(\overline{\mathcal A}=\{H/T:H\in\mathcal A\}\). Thus \(\dim W=3\). The arrangement is called simplicial if every connected component \(C\) of
\[
W\setminus\bigcup_{\overline H\in\overline{\mathcal A}}\overline H
\]
is an open simplicial cone: for some basis \(v_1,v_2,v_3\) of \(W\), depending on \(C\),
\[
C=\{t_1v_1+t_2v_2+t_3v_3:t_1,t_2,t_3>0\}.
\]
\textbf{Principal question.} Are there further rank-three simplicial arrangements, or further infinite series of such arrangements, beyond the examples and series already known in the intended source classification?

To express the first alternative precisely, one must specify an equivalence relation \(\sim\) and a reference collection \(\mathscr K\) of known equivalence classes. With those choices, the question becomes whether there exists a rank-three simplicial arrangement \(\mathcal A\) such that \([\mathcal A]_{\sim}\notin\mathscr K\). Neither choice is supplied. The meaning of a further infinite series also requires clarification: the source does not specify its parameter set, conditions for distinct members, or how novelty relative to known series is determined. Consequently, these missing data prevent a fully specified classification question.

\textbf{Additional target in the original statement.} The original record also asks whether Terao's freeness conjecture holds over a fixed field. This target is absent from the cleaned statement; its intended inclusion remains unresolved. In the conventional finite central setting, fix a field \(K\), let \(S=K[x_1,\ldots,x_\ell]\), and write each hyperplane \(H\) of an arrangement \(\mathcal B\) in \(K^\ell\) as \(H=\ker\alpha_H\), where \(\alpha_H\) is a nonzero linear form. Define
\[
D(\mathcal B)=\{\theta\in\operatorname{Der}_K(S):\theta(\alpha_H)\in\alpha_H S\text{ for every }H\in\mathcal B\}.
\]
The arrangement is free when \(D(\mathcal B)\) is a free \(S\)-module. Its intersection lattice \(L(\mathcal B)\) consists of intersections of subfamilies of \(\mathcal B\), including the empty intersection \(K^\ell\), ordered by reverse inclusion. The corresponding conjectural statement is that, for every fixed field \(K\), every positive integer \(\ell\), and all finite central arrangements \(\mathcal B,\mathcal B'\) in \(K^\ell\),
\[
L(\mathcal B)\cong L(\mathcal B')
\quad\Longrightarrow\quad
\bigl(\mathcal B\text{ is free}\iff\mathcal B'\text{ is free}\bigr).
\]
The supplied wording imposes no characteristic restriction on the fixed field. The central setting used in this formulation, and the real setting used for simpliciality above, are explicit conventional interpretations requiring confirmation rather than additional source hypotheses.
\subsection{Short English statement}
The cleaned question asks whether further rank-three simplicial arrangements or infinite series exist. A precise novelty claim requires the intended domain, equivalence relation, and catalogue of known arrangements. The original statement additionally asks Terao's freeness question over a fixed field; whether that additional target belongs in this entry needs review.
\subsection{Sources}
\begin{itemize}
\item Local UnsolvedMath registry, supplied complete source record: permanent ID \(30004599\), title \emph{Further Rank-Three Simplicial Arrangements}, exact source problem code \texttt{catalogue entry}, in the Oberwolfach Reports Open Problems collection. The record gives proposed year 2021 and no named proposer.
\item Supplied source citation: \emph{Logarithmic Vector Fields and Freeness of Divisors and Arrangements: New perspectives and applications} (2022), Oberwolfach Reports, pp.\ 225--300, DOI \url{https://doi.org/10.4171/owr/2021/5}. The record describes these reports as documenting mathematical workshops and their open-problem discussions. Its literature list labels this source Oberwolfach Report 5/2021; the supplied year labels are retained without reconciliation.
\item Supplied literature reference: Pagaria, \emph{Journal of the London Mathematical Society}, \url{https://doi.org/10.1112/jlms.12829}.
\item Supplied literature reference: Cuntz--Muecksch, \emph{Advances in Applied Mathematics}, \url{https://doi.org/10.1016/j.aam.2019.02.008}.
\item Literature assessment dated exactly \texttt{2026-08-22}: the supplied curation labels the problem open and states that Terao's freeness conjecture remains open even in dimension three and that the classification of rank-three simplicial arrangements is incomplete. This is supplied dated curation derived from the Oberwolfach material, not an independently verified literature assessment; time-sensitive status and cited claims require independent verification.
\end{itemize}
\end{document}
